\renewcommand{\teacherTag}{ANIK SIR \textbullet\ LECTURES 3--4}
\chapter{Convolution and Interconnected Systems}

%------------------------------------------------------------------
\begin{qcard}{2023--2024 \textbullet\ Q1(b)}{CSE 219 \textbullet\ 22 batch p.1}{15}
Consider the combination of discrete-time LTI systems in Figure 1: the input $x[n]$ goes through $h_1$ and $h_2$ in parallel, their outputs are added and passed through $h_3$ to give $y[n]$. The impulse responses are $h_1[n]=\delta[n]-\delta[n-2]$, $h_2[n]=u[n-1]$, $h_3[n]=\delta[n]-\delta[n-1]$. Suppose the input is $x[n]=u[n]$.
(i) Express the equivalent impulse response $h_{eq}[n]$ of the overall system. (ii) Infer the output $y[n]$ for the given input.
\end{qcard}
\begin{solution}
Parallel branches add; the cascade convolves: $h_{eq}=(h_1+h_2)*h_3$.
\begin{align*}
h_2*h_3&=u[n-1]-u[n-2]=\delta[n-1],\\
h_1*h_3&=\delta[n]-\delta[n-1]-\delta[n-2]+\delta[n-3].
\end{align*}
\[h_{eq}[n]=h_1*h_3+h_2*h_3=\delta[n]-\delta[n-2]+\delta[n-3].\]
For $x=u[n]$: $y=u*h_{eq}=u[n]-u[n-2]+u[n-3]$, i.e.
\begin{center}\small\begin{tabular}{c|ccccccc}$n$&$<0$&0&1&2&3&4&$\ge5$\\\hline$y[n]$&0&1&1&0&1&1&1\end{tabular}\end{center}
\ans{$h_{eq}[n]=\delta[n]-\delta[n-2]+\delta[n-3]$; $y[n]=u[n]-u[n-2]+u[n-3]=\{1,1,0,1,1,1,\dots\}$.}
\end{solution}

%------------------------------------------------------------------
\begin{qcard}{2023--2024 \textbullet\ Q1(c)}{EEE 211 \textbullet\ EEE p.1}{17}
Consider the system in the figure: $x(t)$ feeds three parallel paths --- $h_1(t)$; $h_2(t)$ followed by $h_3(t)$; and $h_4(t)$ --- whose outputs are added and passed through $h_5(t)$ to give $y(t)$. Given $h_1=e^{-2t}u(t)$, $h_2=e^{-2t}u(t)$, $h_3=e^{-t}u(t)$, $h_4=\delta(t)$, $h_5=e^{-3t}u(t)$. Using time-domain techniques, (i) determine whether the system is BIBO stable; (ii) determine the output $y(t)$ for the input $x(t)=e^{-t}u(t)$.
\end{qcard}
\begin{solution}
Useful pair: $e^{-at}u*e^{-bt}u=\dfrac{e^{-bt}-e^{-at}}{a-b}u(t)$ for $a\ne b$.

\step{Middle path:} $h_2*h_3=e^{-2t}u*e^{-t}u=(e^{-t}-e^{-2t})u(t)$.
\step{Sum of the three paths:} $g=h_1+h_2*h_3+h_4=e^{-2t}u+(e^{-t}-e^{-2t})u+\delta(t)=\delta(t)+e^{-t}u(t)$.
\step{Overall:} $h_{eq}=g*h_5=e^{-3t}u+e^{-t}u*e^{-3t}u=e^{-3t}u+\tfrac12(e^{-t}-e^{-3t})u$, i.e.
\[h_{eq}(t)=\tfrac12\big(e^{-t}+e^{-3t}\big)u(t).\]
\step{(i)} $h_{eq}$ is causal and $\int_0^\infty\abs{h_{eq}}\dd t=\tfrac12(1+\tfrac13)=\tfrac23<\infty$: BIBO stable.

\step{(ii)} $y=e^{-t}u*h_{eq}=\tfrac12\,e^{-t}u*e^{-t}u+\tfrac12\,e^{-t}u*e^{-3t}u=\tfrac12te^{-t}u+\tfrac14(e^{-t}-e^{-3t})u$.
\ans{(i) Stable ($\int|h_{eq}|=\frac23$). \quad (ii) $y(t)=\big[\tfrac12te^{-t}+\tfrac14e^{-t}-\tfrac14e^{-3t}\big]u(t)$.}
\end{solution}

%------------------------------------------------------------------
\begin{qcard}{2023--2024 \textbullet\ Q2(a)}{EEE 211 \textbullet\ EEE p.2}{18}
Graphically determine the convolution of the pair of signals in the figure: $x(t)=1$ for $1<t<2$ (and 0 elsewhere), and $h(t)$ a triangle of height 1 on $(-1,1)$ with peak at $t=0$.
\end{qcard}
\begin{solution}
Since $x$ is a rectangle on $(1,2)$, $y(t)=\int h(\tau)\,x(t-\tau)\dd\tau=\int_{t-2}^{t-1}h(s)\dd s=H(t-1)-H(t-2)$, where $H(s)=\int_{-\infty}^{s}h$ is the running integral of the triangle:
\[H(s)=\begin{cases}0,&s<-1,\\\frac{(s+1)^2}{2},&-1\le s\le0,\\1-\frac{(1-s)^2}{2},&0\le s\le1,\\1,&s>1.\end{cases}\]
Slide the flipped rectangle (width 1) along $\tau$; the interval table is
\begin{center}\small\begin{tabular}{@{}lll@{}}\toprule interval & overlap & $y(t)$\\\midrule
$t<0$ & none & $0$\\
$0\le t\le1$ & right part of $h$ enters & $\dfrac{t^2}{2}$\\
$1\le t\le2$ & window $(t-2,t-1)$ covers the peak & $-t^2+3t-\dfrac32$\\
$2\le t\le3$ & left part of $h$ leaves & $\dfrac{(3-t)^2}{2}$\\
$t>3$ & none & $0$\\\bottomrule\end{tabular}\end{center}
Checks: continuity at $t=1$ ($\frac12$), $t=2$ ($\frac12$), peak $y(1.5)=\frac34$, and area $\int y=(1)(1)=1$.
\begin{center}
\begin{tikzpicture}\begin{axis}[sig,width=8cm,height=3.6cm,title={$y(t)=x*h$},xmin=-0.5,xmax=3.5,ymin=-0.1,ymax=1,xtick={1,2,3}]
\addplot[thick,navy,domain=0:1]{x^2/2};\addplot[thick,navy,domain=1:2]{-x^2+3*x-1.5};\addplot[thick,navy,domain=2:3]{(3-x)^2/2};
\end{axis}\end{tikzpicture}\end{center}
\ans{$y=\frac{t^2}2$ ($0\le t\le1$), $-t^2+3t-\frac32$ ($1\le t\le2$), $\frac{(3-t)^2}{2}$ ($2\le t\le3$), $0$ otherwise.}
\end{solution}

%------------------------------------------------------------------
\begin{qcard}{2022--2023 \textbullet\ Q8(a)}{CSE 219 \textbullet\ TF p.5}{20}
Convolve $x[n]=3^nu[-n-1]+\big(\tfrac13\big)^nu[n]$ with $h[n]=\big(\tfrac14\big)^nu[n+3]$.
\end{qcard}
\begin{solution}
Split the two-sided input (convolution distributes): $x=x_1+x_2$, $x_1=3^nu[-n-1]$, $x_2=(\frac13)^nu[n]$. The kernel $h[n-k]=(\tfrac14)^{n-k}$ is nonzero for $k\le n+3$.

\step{$y_1=x_1*h$.} $x_1[k]\ne0$ for $k\le-1$, so $k\le\min(-1,n+3)$ and $y_1[n]=4^{-n}\sum_k12^k$:
\begin{itemize}[leftmargin=1.4em,itemsep=2pt]
\item $n\le-4$ (upper limit $n+3$): $y_1[n]=4^{-n}\dfrac{12^{\,n+3}}{1-\frac1{12}}=\dfrac{12}{11}\,4^{-n}12^{\,n+3}=\dfrac{12^4}{11}\,3^{n}=\dfrac{256}{11}\,3^{\,n+4}$.
\item $n\ge-3$ (upper limit $-1$): $y_1[n]=4^{-n}\dfrac{12^{-1}}{1-\frac1{12}}=\dfrac{4^{-n}}{11}$.
\end{itemize}

\step{$y_2=x_2*h$.} $x_2[k]\ne0$ for $k\ge0$, so $0\le k\le n+3$ (needs $n\ge-3$):
\[y_2[n]=4^{-n}\sum_{k=0}^{n+3}\Big(\tfrac43\Big)^k=4^{-n}\,\frac{(4/3)^{n+4}-1}{1/3}=256\cdot3^{-(n+3)}-3\cdot4^{-n}.\]
\step{Total.}
\[
y[n]=\begin{cases}\dfrac{256}{11}\,3^{\,n+4},&n\le-4,\\[8pt]256\cdot3^{-(n+3)}-\dfrac{32}{11}\,4^{-n},&n\ge-3.\end{cases}
\]
For $n\ge-3$: $y_1+y_2=\frac{4^{-n}}{11}+256\cdot3^{-(n+3)}-3\cdot4^{-n}$ and $\frac1{11}-3=-\frac{32}{11}$. Values: $y[-4]=23.27$, $y[-3]=69.82$ (peak), $y[-2]=38.79$, $y[-1]=16.81$, $y[0]=6.57$, $y[1]=2.43$.
\begin{center}
\begin{tikzpicture}\begin{axis}[width=10cm,height=4cm,axis lines=middle,xlabel={$n$},ticklabel style={font=\scriptsize},ymin=0,ymax=80,xmin=-9,xmax=6,xtick distance=1,ycomb,mark=*]
\addplot+[ycomb,navy,mark=*,mark size=1.4pt]coordinates{(-8,0.288)(-7,0.864)(-6,2.59)(-5,7.77)(-4,23.27)(-3,69.82)(-2,38.79)(-1,16.81)(0,6.57)(1,2.43)(2,0.872)(3,0.306)};
\end{axis}\end{tikzpicture}\end{center}
\ans{$y[n]=\frac{256}{11}3^{n+4}$ for $n\le-4$; $\;y[n]=256\cdot3^{-(n+3)}-\frac{32}{11}4^{-n}$ for $n\ge-3$.}
\end{solution}

%------------------------------------------------------------------
\begin{qcard}{2022--2023 \textbullet\ Q7(b)}{EEE 211 \textbullet\ EEE p.11}{15}
Consider an LTI system where $x(t)=h(t)=t\,[u(t)-u(t-1)]$. Find the expression of the output $y(t)$ of the system for all $t$.
\end{qcard}
\begin{solution}
Both signals are the ramp $t$ on $(0,1)$. $y(t)=\int x(\tau)h(t-\tau)\dd\tau$ with $0<\tau<1$ and $0<t-\tau<1$, i.e.\ $\max(0,t-1)<\tau<\min(1,t)$.
\begin{itemize}[leftmargin=1.4em,itemsep=2pt]
\item $0\le t\le1$ (limits $0$ to $t$): $y=\int_0^t\tau(t-\tau)\dd\tau=\dfrac{t^3}{6}$.
\item $1\le t\le2$ (limits $t-1$ to $1$): $y=\int_{t-1}^{1}\tau(t-\tau)\dd\tau=-\dfrac{t^3}{6}+t-\dfrac23$.
\item otherwise $0$.
\end{itemize}
Check: both formulas give $\frac16$ at $t=1$; the second gives $0$ at $t=2$; the area of $y$ is $\frac12\cdot\frac12=\frac14$.
\begin{center}\begin{tikzpicture}\begin{axis}[sig,width=8cm,height=3.4cm,title={$y(t)$},xmin=-0.5,xmax=2.5,ymin=-0.03,ymax=0.25,xtick={1,2}]
\addplot[thick,navy,domain=0:1]{x^3/6};\addplot[thick,navy,domain=1:2]{-x^3/6+x-2/3};\end{axis}\end{tikzpicture}\end{center}
\ans{$y(t)=\frac{t^3}{6}$ for $0\le t\le1$; $\;-\frac{t^3}{6}+t-\frac23$ for $1\le t\le2$; $0$ otherwise.}
\end{solution}

%------------------------------------------------------------------
\begin{qcard}{2021--2022 \textbullet\ Q2(a)}{EEE 211 \textbullet\ EEE p.12}{17}
The input $x(t)$ and the impulse response $h(t)$ of an LTI system are shown in the figure: $x(t)$ is a triangle on $(0,2)$ with peak 1 at $t=1$; $h(t)$ is a rectangle of height 2 on $(0,1)$ together with a unit impulse at $t=3$. Compute the output $y(t)$ using the graphical method. Also write the analytical expression of $y(t)$.
\end{qcard}
\begin{solution}
$h(t)=2[u(t)-u(t-1)]+\delta(t-3)$, so $y=2\int_{t-1}^{t}x(s)\dd s+x(t-3)$. The running integral of the triangle is $X(s)=\frac{s^2}{2}$ ($0\le s\le1$), $1-\frac{(2-s)^2}{2}$ ($1\le s\le2$), $1$ ($s\ge2$).
\begin{center}\small\begin{tabular}{@{}ll@{}}\toprule interval & rectangular part $2[X(t)-X(t-1)]$\\\midrule
$0\le t\le1$ & $t^2$\\
$1\le t\le2$ & $2-(2-t)^2-(t-1)^2=-2t^2+6t-3$\\
$2\le t\le3$ & $(3-t)^2$\\\bottomrule\end{tabular}\end{center}
The impulse produces the delayed triangle $x(t-3)$: $t-3$ on $(3,4)$ and $5-t$ on $(4,5)$.
\begin{center}\begin{tikzpicture}\begin{axis}[sig,width=9cm,height=3.6cm,title={$y(t)$},xmin=-0.5,xmax=5.5,ymin=-0.1,ymax=1.7,xtick={1,2,3,4,5}]
\addplot[thick,navy,domain=0:1]{x^2};\addplot[thick,navy,domain=1:2]{-2*x^2+6*x-3};\addplot[thick,navy,domain=2:3]{(3-x)^2};
\addplot[thick,navy]coordinates{(3,0)(4,1)(5,0)};\end{axis}\end{tikzpicture}\end{center}
\ans{$y=t^2$ ($0\le t\le1$); $-2t^2+6t-3$ ($1\le t\le2$); $(3-t)^2$ ($2\le t\le3$); then the triangle $x(t-3)$ on $(3,5)$.}
\end{solution}

%------------------------------------------------------------------
\begin{qcard}{2020--2021 \textbullet\ Q6(b)}{EEE 211 \textbullet\ EEE p.20}{10}
Using graphical convolution, determine the output $y(t)$ of an LTI system if the input is $x(t)=\operatorname{rect}(t/2)$ and the system impulse response is $h(t)=r(t)\,u(-t+1)$. Also write the analytical expression of $y(t)$ and draw it.
\end{qcard}
\begin{solution}
$x=1$ on $(-1,1)$ and $h(t)=t$ on $(0,1)$. With $x$ a rectangle, $y(t)=\int_{t-1}^{t+1}h(s)\dd s$.
\begin{center}\small\begin{tabular}{@{}lll@{}}\toprule interval & overlap of $(t-1,t+1)$ with $(0,1)$ & $y(t)$\\\midrule
$-1\le t\le0$ & $(0,t+1)$ & $\frac{(t+1)^2}{2}$\\
$0\le t\le1$ & the whole $(0,1)$ & $\frac12$\\
$1\le t\le2$ & $(t-1,1)$ & $\frac12-\frac{(t-1)^2}{2}$\\
else & none & $0$\\\bottomrule\end{tabular}\end{center}
\begin{center}\begin{tikzpicture}\begin{axis}[sig,width=8cm,height=3.2cm,title={$y(t)$},xmin=-2,xmax=3,ymin=-0.1,ymax=0.7,xtick={-1,1,2}]
\addplot[thick,navy,domain=-1:0]{(x+1)^2/2};\addplot[thick,navy,domain=0:1]{0.5};\addplot[thick,navy,domain=1:2]{0.5-(x-1)^2/2};\end{axis}\end{tikzpicture}\end{center}
\ans{$y=\frac{(t+1)^2}{2}$ ($-1\le t\le0$), $\frac12$ ($0\le t\le1$), $\frac{1-(t-1)^2}{2}$ ($1\le t\le2$), $0$ otherwise.}
\end{solution}

%------------------------------------------------------------------
\begin{qcard}{2019--2020 \textbullet\ Q2}{EEE 211 \textbullet\ EEE p.22--23}{35}
Suppose an LTI system has impulse response $h(t)$ and input $x(t)$. The figure shows $x(t)$ (a rectangle of height 1 on $(0,2)$) and $h(t+1)$ (a triangle of height 1 on $(-1,1)$).
(i) Is the system causal? (ii) Use the graphical convolution technique to determine $y(t)=x(t)*h(t)$. (iii) Compare $x(t)$ and $y(t)$, and explain what this system does. (iv) Sketch the output of the system for a different input $z(t)=x(t)-x(t-4)$.
\end{qcard}
\begin{solution}
\step{(i)} $h(t+1)$ is a triangle on $(-1,1)$, so $h(t)$ is a triangle on $(0,2)$ with peak at $t=1$: $h(t)=0$ for $t<0$ --- \emph{causal}.

\step{(ii)} $x$ is a rectangle on $(0,2)$: $y(t)=\int_{t-2}^{t}h(s)\dd s=H(t)-H(t-2)$ with $H(s)=\frac{s^2}2$ ($0\le s\le1$), $1-\frac{(2-s)^2}{2}$ ($1\le s\le2$), $1$ ($s\ge2$):
\[
y(t)=\begin{cases}\frac{t^2}2,&0\le t\le1,\\1-\frac{(2-t)^2}2,&1\le t\le2,\\1-\frac{(t-2)^2}2,&2\le t\le3,\\\frac{(4-t)^2}2,&3\le t\le4,\\0,&\text{otherwise.}\end{cases}
\]
\step{(iii)} $y$ is a smoothed copy of $x$: the sharp edges become S-shaped curves, the duration grows from 2 to 4 (sum of the durations), the peak is $y(2)=1$ and the area is preserved ($\int y=\int x\int h=2\cdot1=2$). The system is a causal smoothing (low-pass) filter: a triangular-weighted moving average over 2 seconds with a delay of 1.

\step{(iv)} $z=x(t)-x(t-4)$, so the output is $y(t)-y(t-4)$: the smooth pulse on $(0,4)$ followed by its negative on $(4,8)$.
\begin{center}
\begin{tikzpicture}\begin{axis}[sig,width=7.5cm,height=3.2cm,title={$y(t)=x*h$},xmin=-0.5,xmax=4.5,ymin=-0.1,ymax=1.3,xtick={1,2,3,4}]
\addplot[thick,navy,domain=0:1]{x^2/2};\addplot[thick,navy,domain=1:2]{1-(2-x)^2/2};\addplot[thick,navy,domain=2:3]{1-(x-2)^2/2};\addplot[thick,navy,domain=3:4]{(4-x)^2/2};\end{axis}\end{tikzpicture}\hspace{2mm}
\begin{tikzpicture}\begin{axis}[sig,width=7.5cm,height=3.2cm,title={output for $z(t)$},xmin=-0.5,xmax=8.5,ymin=-1.3,ymax=1.3,xtick={2,4,6,8}]
\addplot[thick,navy,domain=0:1]{x^2/2};\addplot[thick,navy,domain=1:2]{1-(2-x)^2/2};\addplot[thick,navy,domain=2:3]{1-(x-2)^2/2};\addplot[thick,navy,domain=3:4]{(4-x)^2/2};
\addplot[thick,crimson,domain=4:5]{-(x-4)^2/2};\addplot[thick,crimson,domain=5:6]{-(1-(2-(x-4))^2/2)};\addplot[thick,crimson,domain=6:7]{-(1-((x-4)-2)^2/2)};\addplot[thick,crimson,domain=7:8]{-(8-x)^2/2};\end{axis}\end{tikzpicture}
\end{center}
\ans{Causal; $y$ is a 4-second smooth pulse (peak 1 at $t=2$), i.e.\ a smoothing low-pass system; for $z$ the output is $y(t)-y(t-4)$.}
\end{solution}

%------------------------------------------------------------------
\begin{qcard}{2018--2019 \textbullet\ Q7(b)}{EEE 211 \textbullet\ EEE p.31}{15}
Let $h(t)=u(t)-u(t-1)$ and $x(t)=t^2\,[u(t)-u(t-3)]$. Using graphical convolution compute the output $y(t)$ and write its analytical expression.
\end{qcard}
\begin{solution}
$h$ is a rectangle of width 1, so $y(t)=\int_{t-1}^{t}x(s)\dd s$ with $x=s^2$ on $(0,3)$:
\begin{center}\small\begin{tabular}{@{}lll@{}}\toprule interval & window $(t-1,t)$ & $y(t)$\\\midrule
$0\le t\le1$ & $(0,t)$ & $\dfrac{t^3}{3}$\\
$1\le t\le3$ & inside $(0,3)$ & $\dfrac{t^3-(t-1)^3}{3}=t^2-t+\dfrac13$\\
$3\le t\le4$ & $(t-1,3)$ & $\dfrac{27-(t-1)^3}{3}$\\
else & none & $0$\\\bottomrule\end{tabular}\end{center}
Check: continuity at $t=1$ ($\frac13$) and $t=3$ ($\frac{19}{3}$), $y(4)=0$, area $=9\cdot1=9$.
\begin{center}\begin{tikzpicture}\begin{axis}[sig,width=8cm,height=3.4cm,title={$y(t)$},xmin=-0.5,xmax=4.5,ymin=-0.5,ymax=8,xtick={1,3,4}]
\addplot[thick,navy,domain=0:1]{x^3/3};\addplot[thick,navy,domain=1:3]{x^2-x+1/3};\addplot[thick,navy,domain=3:4]{(27-(x-1)^3)/3};\end{axis}\end{tikzpicture}\end{center}
\ans{$y=\frac{t^3}3$ ($0\le t\le1$); $t^2-t+\frac13$ ($1\le t\le3$); $\frac{27-(t-1)^3}{3}$ ($3\le t\le4$); $0$ otherwise.}
\end{solution}

%------------------------------------------------------------------
\begin{qcard}{2018--2019 \textbullet\ Q4(b)}{CSE 311 \textbullet\ TF p.18}{17}
Prove that for an input signal $x(t)$, the output of a linear time-invariant system with impulse response $h(t)$ is given by $y(t)=x(t)*h(t)$. Use \emph{only} time-domain variables and equations.
\end{qcard}
\begin{solution}
Let $S$ denote the system and $h(t)=S\{\delta(t)\}$ (definition of the impulse response).
\begin{steps}
\item \emph{Sifting.} Any input is a continuum of impulses: $x(t)=\displaystyle\int_{-\infty}^{\infty}x(\tau)\,\delta(t-\tau)\dd\tau$.
\item \emph{Time invariance.} $S\{\delta(t-\tau)\}=h(t-\tau)$.
\item \emph{Homogeneity.} The weight $x(\tau)\dd\tau$ does not depend on $t$, so $S\{x(\tau)\dd\tau\,\delta(t-\tau)\}=x(\tau)\dd\tau\,h(t-\tau)$.
\item \emph{Additivity (the integral is a limit of sums).}
\[y(t)=S\{x(t)\}=\int_{-\infty}^{\infty}x(\tau)\,h(t-\tau)\dd\tau=(x*h)(t).\]
\end{steps}
Only linearity (steps 3--4) and time invariance (step 2) were used, so the result holds for every LTI system. (Discrete time: $x[n]=\sum_kx[k]\delta[n-k]\Rightarrow y[n]=\sum_kx[k]h[n-k]$.)
\ans{$y(t)=\int x(\tau)h(t-\tau)\dd\tau=x*h$, from sifting, time invariance and linearity.}
\end{solution}

